% ============================================================
%  第3章：对偶理论
% ============================================================
\section{第3章 \quad 对偶理论}

\subsection{对偶问题的构造}

\textbf{对称形式（$\min$ 原问题，$\le$ 约束）：}
\begin{center}
\begin{tabular}{c|c}
\textbf{原问题 (P)} & \textbf{对偶问题 (D)} \\
\hline
$\min\ c^T x$ & $\max\ b^T y$ \\
$\text{s.t.}\ Ax \ge b$ & $\text{s.t.}\ A^T y \le c$ \\
$x \ge 0$ & $y \ge 0$ \\
\end{tabular}
\end{center}

\textbf{一般构造规则：}
\begin{itemize}
  \item 原问题 $\min$ $\to$ 对偶问题 $\max$
  \item 原问题第 $i$ 约束的右端 $b_i$ $\to$ 对偶变量 $y_i$
  \item 原问题第 $j$ 变量的系数 $c_j$ $\to$ 对偶约束的右端
  \item 等式约束 $\to$ 对偶变量无符号限制
  \item $\ge$ 约束（对 $\min$ 问题）$\to$ 对偶变量 $\le 0$
  \item 自由变量 $\to$ 对偶等式约束
\end{itemize}

\subsection{对偶定理}

\textbf{定理 3.1（弱对偶定理）}
设 $x$ 和 $y$ 分别是 (P) 和 (D) 的可行解，则 $c^T x \ge b^T y$。
推论：若 $c^T x = b^T y$，则 $x$ 和 $y$ 分别是各自问题的最优解。

\textbf{定理 3.2（强对偶定理）}
若 (P) 有有限最优解 $x^*$，则 (D) 也有最优解 $y^*$，且 $c^T x^* = b^T y^*$。

\subsection{互补松弛条件}

\textbf{定理 3.3（互补松弛）}
设 $x^*$ 和 $y^*$ 分别是 (P) 和 (D) 的可行解。它们同时为最优的充要条件是：
\[
(y^*)^T(Ax^* - b) = 0,\qquad (c - A^T y^*)^T x^* = 0
\]
即：
\begin{itemize}
  \item 若 $x_j^* > 0$，则对应对偶约束取等号
  \item 若对偶约束严格不取等，则 $x_j^* = 0$
  \item 若 $y_i^* > 0$，则对应原约束取等号
  \item 若原约束严格不取等，则 $y_i^* = 0$
\end{itemize}

\subsection{对偶单纯形法}

\textbf{适用场景：} 原始不可行但对偶可行（检验数全 $\ge 0$ 但 $B^{-1}b$ 有负分量）。

\textbf{算法步骤：}
\begin{enumerate}
  \item 确定出基变量：选 $B^{-1}b$ 中最负的分量对应行
  \item 确定进基变量：在该行负系数列中，取 $\displaystyle \min\left\{\frac{\sigma_j}{|a_{rj}|} \;\middle|\; a_{rj} < 0\right\}$
  \item 转轴运算（同单纯形法）
  \item 重复直至 $B^{-1}b \ge 0$（原始可行）
\end{enumerate}

\subsection{影子价格}

\textbf{定义：} 对偶变量 $y_i^*$ 的经济含义——第 $i$ 种资源增加 1 单位时，最优目标值的边际变化量。$y_i^* = \frac{\partial z^*}{\partial b_i}$。

\subsection{本章速查}

\begin{itemize}
  \item \textbf{对称对偶}：$\min\ c^Tx,\ Ax\ge b,\ x\ge0 \;\longleftrightarrow\; \max\ b^Ty,\ A^Ty\le c,\ y\ge0$
  \item \textbf{弱对偶}：$c^Tx \ge b^Ty$；\textbf{强对偶}：最优值相等
  \item \textbf{互补松弛}：$(y^*)^T(Ax^*-b)=0$，$(c-A^Ty^*)^T x^*=0$
  \item \textbf{对偶单纯形法}：原始不可行但对偶可行时使用
  \item \textbf{影子价格}：$y_i^* = \partial z^* / \partial b_i$
\end{itemize}
